% --- Styling for the prompt boxes ------------------------------------------
\definecolor{promptblue}{HTML}{EAF2FF}
\definecolor{promptgreen}{HTML}{ECF8F1}
\definecolor{promptorange}{HTML}{FFF4E8}
\definecolor{frameblue}{HTML}{2E5AAC}
\definecolor{framegreen}{HTML}{2F7D4A}
\definecolor{frameorange}{HTML}{B86A00}
\definecolor{fragmentgray}{HTML}{F7F7F8}
\definecolor{fragmentframe}{HTML}{8A8F98}

\tcbset{
  enhanced,
  breakable,
  boxrule=0.7pt,
  arc=1.8mm,
  left=2mm,right=2mm,top=1.5mm,bottom=1.5mm,
  coltitle=white,
  fonttitle=\bfseries,
}

\newtcolorbox{stageonebox}{%
  colback=promptblue,
  colframe=frameblue,
  title={Stage 1 Prompt: Reflection},
}
\newtcolorbox{stagetwobox}{%
  colback=promptgreen,
  colframe=framegreen,
  title={Stage 2 Prompt: Belief Update},
}
\newtcolorbox{stagethreebox}{%
  colback=promptorange,
  colframe=frameorange,
  title={Stage 3 Prompt: Decision},
}
\newtcolorbox{fragmentbox}[1]{%
  colback=fragmentgray,
  colframe=fragmentframe,
  title={#1},
}

\newcommand{\promptfont}{\ttfamily\small}
\newcommand{\ph}[1]{\textit{\{#1\}}}


% ===========================================================================
\section{Ground Truth Model Construction}
\label{app:armote}
% ===========================================================================

\subsection{The ARMOTE-MO Workflow}
\label{app:armote_algorithm}

\begin{algorithm}[H]
\caption{ARMOTE-MO: automated multi-output regression with multi-objective
         hyperparameter optimization}
\label{alg:armote}
\begin{algorithmic}[1]
\Require Composition--property dataset $\mathcal{D} = \{(\mathbf{x}_i, \mathbf{y}_i)\}_{i=1}^{N}$,
         search space $\Theta$, trial budget $T = 100$
\State Split $\mathcal{D}$ into $\mathcal{D}_\mathrm{train}$ / $\mathcal{D}_\mathrm{test}$ (80/20)
\State Fit standardization scalers for inputs and outputs on $\mathcal{D}_\mathrm{train}$;
       apply them to $\mathcal{D}_\mathrm{train}$ and $\mathcal{D}_\mathrm{test}$
\For{$t = 1$ to $T$}
    \State Sample $\theta_t \in \Theta$ with the Tree-Parzen Estimator
    \State Evaluate $\theta_t$ by 5-fold cross-validation on the standardized
           $\mathcal{D}_\mathrm{train}$, recording the mean MSE (standardized targets)
           and the mean $R^2$, each averaged across outputs
\EndFor
\State $\Theta^{*} \gets$ non-dominated set of $\{\theta_t\}$ in (MSE, $R^2$) space
       (minimize MSE, maximize $R^2$)
\State $\theta^{\dagger} \gets \arg\max_{\theta \in \Theta^{*}}$ cross-validated mean $R^2$
       (averaged across outputs)
\State Refit the multi-output RFR on all of the standardized $\mathcal{D}_\mathrm{train}$ at $\theta^{\dagger}$
\State \Return model, fitted scalers, and held-out $R^2$ on $\mathcal{D}_\mathrm{test}$
       (per output and averaged across outputs, in original target units)
\end{algorithmic}
\end{algorithm}

\subsection{Hyperparameter Search Spaces}
\label{app:armote_spaces}

\begin{table}[H]
\centering
\begin{tabular}{llc}
\toprule
Hyperparameter & Search range & Scale \\
\midrule
$n_\mathrm{trees}$           & [10, 150] & integer, uniform \\
\texttt{max\_depth}          & [10, 150] & integer, uniform \\
\texttt{min\_samples\_split} & [2, 20]   & integer, uniform \\
\texttt{min\_samples\_leaf}  & [1, 20]   & integer, uniform \\
\bottomrule
\end{tabular}
\caption{Hyperparameter search spaces explored by ARMOTE-MO for both alloy
         systems. Ranges are inclusive. All other random forest settings,
         including \texttt{max\_features}, were left at their scikit-learn
         defaults. The selected configurations are reported in
         Table~\ref{tab:rfr-hyperparameters}.}
\label{tab:armote-search-space}
\end{table}

\subsection{Truth-Function Generated Datasets}
\label{app:truth_datasets}

\section{Prompts Used by the LLM Allocation Agents}
\label{app:prompts}

\noindent The two LLM-based allocation strategies used separate two-call prompt
pipelines. Dynamic quantities inserted by the campaign code are shown in braces.
All prompt calls used GPT-4o through LangChain. The RHEA campaigns used
temperature $T = 0.2$, and the Fe-Co-Ni magnet campaigns used temperature
$T = 0.3$.

\subsection{Adaptive-\texorpdfstring{$\beta$}{β} Agent Prompts}
\label{app:prompts_adaptive_beta}

\begin{fragmentbox}{Adaptive-$\beta$ Call 1: $\beta$ selection}
\promptfont
You are selecting a beta parameter for a Bayesian Optimization campaign.
\vspace{0.35em}
Problem: \ph{problem\_description}\\
Objectives: \ph{objectives\_str}
\vspace{0.35em}
qEHVI focuses on what the GP currently believes is near-Pareto-optimal. It is the
reliable Pareto front improver and the right default in most situations.
\vspace{0.35em}
qUCB with elevated $\beta$ scouts design-space regions the GP has not characterized
well --- regions with high posterior uncertainty. Its purpose is to find productive
territory qEHVI has not yet been directed toward. Exploration will often produce no
short-term HV gain. That is acceptable when it locates new productive regions faster
than qEHVI alone would. It is NOT acceptable when exploration is saturated (MI
declining) or when the campaign is nearly complete.
\vspace{0.35em}
$\beta$ controls how far into uncertain regions qUCB will probe, measured in
posterior standard deviations. GP inputs and outputs are normalized, so $\beta$ has
the same meaning regardless of problem scale or physical units.
\vspace{0.35em}
\hspace*{1em}$\beta \approx 2$: Conservative. qUCB candidates stay near regions the GP already knows well. Right default when HV is actively improving.\\
\hspace*{1em}$\beta = 3$--$5$: Moderate exploration. Noticeably more uncertainty-seeking. Appropriate when HV progress has slowed but MI signals remaining territory.\\
\hspace*{1em}$\beta = 5$--$8$: Aggressive scouting. Sends candidates into poorly-characterized regions. Justified when HV has materially stalled and MI is still elevated.\\
\hspace*{1em}$\beta > 8$: Only in extreme stagnation with a strong MI signal.
\vspace{0.35em}
Early in the campaign (first \~{}20--25\% of resources consumed), the GP posterior is
poorly constrained --- uncertainty estimates are unreliable and elevated beta may send
candidates to arbitrary regions rather than genuinely informative ones. Prefer
$\beta$ near 2 unless you observe a specific signal that justifies deviation.
\vspace{0.35em}
$\beta$ only has effect if the selected option also allocates exploration points.
Choosing high $\beta$ with 0 explore points in the batch is a no-op.
\vspace{0.35em}
=== CAMPAIGN STATE ===\\
Iteration: \ph{state.iteration}\\
Resources consumed: \ph{budget\_pct}\%  |  Time remaining: \ph{time\_pct\_str}  |  Iters remaining: \~{}\ph{iters\_remaining}\\
Pareto front HV: \ph{hv\_line}\\
Pareto size: \ph{pareto\_size} points (\ph{pareto\_change})
\vspace{0.35em}
=== OPTIMIZATION SIGNALS ===\\
HV velocity ratio: \ph{vel\_str}\\
\hspace*{1em}(recent mean $\Delta$HV divided by historical baseline --- 1.0 = on pace, <0.5 = materially slowed, >1.0 = accelerating)\\
MI of best explore option: \ph{mi\_line} (\ph{mi\_vs\_str})\\
\hspace*{1em}(declining MI = design space increasingly well-characterized; elevated MI = unexplored territory remains)\\
\ph{pareto\_coverage}\\
=== RECENT HISTORY ===\\
\ph{recent\_history}
\vspace{0.35em}
=== OUTPUT FORMAT ===\\
Output EXACTLY these two lines (nothing else):\\
BETA: <number>\\
REASONING: <one sentence citing specific signals above>
\end{fragmentbox}

\begin{fragmentbox}{Adaptive-$\beta$ Call 2: allocation option selection}
\promptfont
You are selecting a resource allocation option for a Bayesian Optimization campaign.
\vspace{0.35em}
=== CONTEXT ===\\
Beta selected: \ph{beta:.1f}\\
Your reasoning: \ph{call1\_reasoning}\\
Resources consumed: \ph{budget\_pct}\%  |  Time remaining: \ph{time\_pct\_str}\\
Iters remaining: \~{}\ph{iters\_remaining}\\
Pareto front: \ph{hv\_pts\_str} pts, HV=\ph{hv\_hv\_str}
\vspace{0.35em}
=== HOW TO READ THE OPTIONS TABLE ===\\
Each option allocates the batch between qEHVI (exploitation) and qUCB (exploration) points.\\
qEHVI score: expected direct Pareto HV improvement --- compare qEHVI values across options only.\\
MI score: expected information gain about unexplored regions --- compare MI values across options only.\\
Do NOT compare qEHVI against MI --- they are on different scales.\\
Higher-index options send more points to qUCB at $\beta$=\ph{beta:.1f}.\\
Options with more qUCB points trade qEHVI improvement for MI gain.
\vspace{0.35em}
\ph{options\_table}
\vspace{0.35em}
=== OUTPUT FORMAT ===\\
Output EXACTLY these two lines (nothing else):\\
SELECTED\_OPTION: <integer 0 to \ph{num\_options - 1}>\\
REASONING: <1-2 sentences on why this option's qEHVI/MI tradeoff fits your beta reasoning>
\end{fragmentbox}

\subsection{Rule-Swap Agent Prompts}
\label{app:prompts_rule_swap}

\begin{fragmentbox}{Rule-Swap Call 1: assessment and $\beta$ selection}
\promptfont
You are guiding a Bayesian Optimization campaign.\\
Problem: \ph{problem\_description}\\
Objectives: \ph{objectives\_str}
\vspace{0.35em}
qEHVI already balances exploration and exploitation within the regions it believes the
Pareto front is likely to expand into. It naturally pushes toward unexplored areas
when those areas have high improvement potential. It is the right default in most
situations.
\vspace{0.35em}
qUCB serves a different purpose: it explores regions of the design space that have not
been evaluated much, regardless of whether the GP currently predicts them to be
Pareto-improving. Its value is in discovering areas the model has not seen yet ---
regions that may contain Pareto-optimal points qEHVI has never been directed toward.
\vspace{0.35em}
$\beta$ controls how explorative qUCB is: larger $\beta$ weights selection toward
points with greater GP uncertainty, pushing further into unmapped territory. Default
$\beta = 2.0$.
\vspace{0.35em}
MutualInfo measures how much the batch reduces GP uncertainty about the optimum. The
absolute value is not interpretable --- only compare it across options within the same
iteration. A higher MutualInfo option queries regions the GP understands least,
regions qEHVI tends to skip because it follows the surrogate's existing beliefs. It
does not directly guarantee better Pareto results, but it signals where the GP's
knowledge is thinnest and where qEHVI is most likely to be under-sampling.
\vspace{0.35em}
=== ITERATION HISTORY (last \ph{n\_completed} completed) ===\\
\ph{history\_text}
\vspace{0.35em}
=== RESOURCES ===\\
Iterations completed: \ph{state.iteration}\\
Estimated iters left: \ph{iters\_left}\\
Budget remaining: \$\ph{state.budget\_remaining:.2f} (\$\ph{state.cost\_per\_point:.2f}/point)\\
Time remaining: \ph{state.time\_remaining:.1f} weeks\\
Current Pareto front: \ph{pareto\_summary}
\vspace{0.35em}
As fewer iterations remain, decisions should shift from probing uncertain regions
toward acting on what the GP already knows is good: \~{}20+ iterations left allows
low-cost exploration, \~{}10 iterations left should prioritize high-confidence
improvements, and \~{}5 or fewer should exploit near-exclusively. If an event reduced
remaining iterations sharply, use the new iteration count.
\vspace{0.35em}
=== YOUR TASK ===\\
Examine the HV/pt column in the history above, but read it carefully in context.
\vspace{0.35em}
STEP 1 --- Identify which iterations were deliberate exploration. Rows where $\beta >
2.0$ or qUCB points > 0 represent intentional exploration investments. Low or zero
HV/pt on those rows is EXPECTED and NORMAL; do not count these rows as evidence of
crowding or stagnation.
\vspace{0.35em}
STEP 2 --- Assess the trend on qEHVI-primary iterations only. Crowding means HV/pt is
declining while Pareto points are still being added; stagnation means zero new Pareto
points across multiple consecutive qEHVI-primary iterations; healthy expansion means
HV/pt is stable or rising.
\vspace{0.35em}
STEP 3 --- Check whether past exploration paid off with a one-iteration lag. If the
most recent iteration used high $\beta$ or qUCB, do NOT raise $\beta$ further yet; give
qEHVI one iteration to exploit the updated model unless the row after exploration also
shows continued stagnation or crowding.
\vspace{0.35em}
STEP 4 --- Factor in iterations remaining. With very few iterations left, raising
$\beta$ is riskier. When in doubt with few iters left, prefer $\beta = 2.0$.
\vspace{0.35em}
Respond using EXACTLY this format:\\
ASSESSMENT: <1-2 sentences describing what the HV/pt trend shows --- crowding, stagnation, or healthy expansion>\\
BETA: <number 1--100; default is 2.0>\\
REASONING: <1 sentence explaining the beta choice>
\end{fragmentbox}

\begin{fragmentbox}{Rule-Swap Call 2: allocation option selection}
\promptfont
You are selecting a resource allocation for a Bayesian Optimization campaign.
\vspace{0.35em}
=== SITUATION ASSESSMENT ===\\
\ph{assessment}
\vspace{0.35em}
=== AVAILABLE OPTIONS ===\\
\ph{options\_text}
\vspace{0.35em}
Note: qEHVI score and MutualInfo are on different scales --- compare each metric
across options separately; do not compare them to each other. The MutualInfo absolute
value is not meaningful on its own: only use it to rank options against each other. A
notably higher MutualInfo on explore-heavy options signals those options query regions
the GP knows least --- use this alongside the crowding/stagnation assessment when
deciding how many qUCB points to include.
\vspace{0.35em}
=== DECISION RULES ===\\
Option 0 (all qEHVI) is the DEFAULT. Choose it unless the assessment above identifies
crowding or stagnation that warrants qUCB exploration.\\
\hspace*{1em}$\bullet$ Crowding detected: choose an option with some qUCB points; scale the number of qUCB points to the severity of the signal.\\
\hspace*{1em}$\bullet$ Stagnation detected: use a larger qUCB share.\\
\hspace*{1em}$\bullet$ Estimated iters left is very low (\~{}5 or fewer): prefer Option 0 regardless of crowding.\\
\hspace*{1em}$\bullet$ No clear signal, or early in the campaign: Option 0.
\vspace{0.35em}
Resources: est. iters left = \ph{iters\_left} | budget \$\ph{state.budget\_remaining:.2f}
\vspace{0.35em}
Respond using EXACTLY this format:\\
SELECTED\_OPTION: <integer 0 to \ph{num\_options - 1}>\\
REASONING: <1-2 sentences connecting the assessment to this specific choice>
\end{fragmentbox}

\subsection{Event Handling in Agent Prompts}
\label{app:prompts_events}

Resource events enter the agent prompts through the campaign-state and remaining-iterations fields above. The Rule-Swap prompt additionally instructs the model to treat an event-shortened campaign according to the new estimated iteration count rather than the original budget.


% ===========================================================================
\section{Supplementary Results}
\label{app:supplementary_results}
% ===========================================================================

\subsection{Per-Experiment Cumulative Mutual Information}
\label{app:cumulative_mi_table}

\begin{table}[H]
\centering
\small
\begin{tabular}{lllccc}
\toprule
System & Experiment & Strategy & Cumulative MI (nats) & Final HV & AUC-HV \\
\midrule
RHEA & Baseline & Adaptive-$\beta$ & $19.54 \pm 0.41$ & $0.9967 \pm 0.00019$ & $51.59 \pm 0.092$ \\
 &  & Rule-Swap & $17.20 \pm 0.50$ & $0.9970 \pm 0.00018$ & $51.79 \pm 0.097$ \\
 &  & Fixed Mixed & $15.48 \pm 0.34$ & $0.9971 \pm 0.00017$ & $50.98 \pm 0.093$ \\
 &  & qEHVI & $14.23 \pm 0.45$ & $0.9983 \pm 0.00018$ & $52.02 \pm 0.077$ \\
\midrule
RHEA & Budget cut & Adaptive-$\beta$ & $12.29 \pm 0.38$ & $0.9843 \pm 0.00071$ & $21.87 \pm 0.078$ \\
 &  & Rule-Swap & $10.93 \pm 0.43$ & $0.9861 \pm 0.00048$ & $21.99 \pm 0.086$ \\
 &  & Fixed Mixed & $11.48 \pm 0.42$ & $0.9855 \pm 0.00043$ & $21.26 \pm 0.070$ \\
 &  & qEHVI & $8.78 \pm 0.31$ & $0.9869 \pm 0.00063$ & $22.11 \pm 0.084$ \\
\midrule
RHEA & Time cut & Adaptive-$\beta$ & $11.97 \pm 0.37$ & $0.9848 \pm 0.00054$ & $21.84 \pm 0.081$ \\
 &  & Rule-Swap & $10.53 \pm 0.35$ & $0.9858 \pm 0.00044$ & $21.95 \pm 0.062$ \\
 &  & Fixed Mixed & $10.63 \pm 0.28$ & $0.9846 \pm 0.00054$ & $21.30 \pm 0.093$ \\
 &  & qEHVI & $8.86 \pm 0.39$ & $0.9866 \pm 0.00100$ & $22.00 \pm 0.089$ \\
\midrule
RHEA & Dual & Adaptive-$\beta$ & $6.11 \pm 0.35$ & $0.9341 \pm 0.00602$ & $7.33 \pm 0.073$ \\
 &  & Rule-Swap & $5.25 \pm 0.26$ & $0.9413 \pm 0.00423$ & $7.37 \pm 0.073$ \\
 &  & Fixed Mixed & $5.32 \pm 0.18$ & $0.9488 \pm 0.00216$ & $7.10 \pm 0.072$ \\
 &  & qEHVI & $4.66 \pm 0.33$ & $0.9486 \pm 0.00318$ & $7.43 \pm 0.075$ \\
\midrule
Fe-Co-Ni & Baseline & Adaptive-$\beta$ & $19.63 \pm 0.48$ & $0.8539 \pm 0.00976$ & $38.23 \pm 0.317$ \\
 &  & Rule-Swap & $16.74 \pm 0.47$ & $0.8457 \pm 0.00871$ & $38.26 \pm 0.254$ \\
 &  & Fixed Mixed & $15.34 \pm 0.48$ & $0.8555 \pm 0.00963$ & $38.09 \pm 0.288$ \\
 &  & qEHVI & $13.17 \pm 0.34$ & $0.8438 \pm 0.00932$ & $38.37 \pm 0.306$ \\
\midrule
Fe-Co-Ni & Budget cut & Adaptive-$\beta$ & $14.60 \pm 0.50$ & $0.8186 \pm 0.00880$ & $18.26 \pm 0.122$ \\
 &  & Rule-Swap & $11.99 \pm 0.47$ & $0.8057 \pm 0.00900$ & $18.20 \pm 0.203$ \\
 &  & Fixed Mixed & $10.82 \pm 0.40$ & $0.7992 \pm 0.00627$ & $18.05 \pm 0.132$ \\
 &  & qEHVI & $9.67 \pm 0.37$ & $0.8155 \pm 0.00850$ & $18.46 \pm 0.180$ \\
\midrule
Fe-Co-Ni & Time cut & Adaptive-$\beta$ & $14.66 \pm 0.50$ & $0.8018 \pm 0.00618$ & $17.37 \pm 0.095$ \\
 &  & Rule-Swap & $12.45 \pm 0.58$ & $0.8178 \pm 0.00874$ & $17.62 \pm 0.123$ \\
 &  & Fixed Mixed & $11.43 \pm 0.62$ & $0.7941 \pm 0.00547$ & $17.17 \pm 0.156$ \\
 &  & qEHVI & $10.44 \pm 0.48$ & $0.7999 \pm 0.00721$ & $17.44 \pm 0.131$ \\
\midrule
Fe-Co-Ni & Dual & Adaptive-$\beta$ & $14.69 \pm 0.50$ & $0.8033 \pm 0.00751$ & $18.11 \pm 0.136$ \\
 &  & Rule-Swap & $12.28 \pm 0.34$ & $0.8111 \pm 0.00854$ & $18.41 \pm 0.163$ \\
 &  & Fixed Mixed & $10.83 \pm 0.42$ & $0.8046 \pm 0.00697$ & $18.08 \pm 0.185$ \\
 &  & qEHVI & $9.97 \pm 0.39$ & $0.8056 \pm 0.00697$ & $18.39 \pm 0.143$ \\
\bottomrule
\end{tabular}

\caption{Cumulative mutual information, final normalized hypervolume, and area
         under the hypervolume curve for every system, experiment, and strategy
         (mean $\pm$ SEM, $n = 25$ seeds).}
\label{tab:cumulative-mi-full}
\end{table}

\subsection{Pairwise Significance Tests}
\label{app:significance}

All strategy comparisons use the two-sided Wilcoxon signed-rank test applied to
per-seed differences. The two strategies in a pair are matched by shared random seed,
so each of the $n = 25$ paired observations contrasts strategies that began from the
same initial design and the same random state. No multiple-comparison correction is
applied: the 144 tests reported in Table~\ref{tab:significance} (six strategy pairs
$\times$ three metrics $\times$ eight experiments) are individually uncorrected, and
$p$-values close to the conventional $0.05$ threshold should therefore be read as
descriptive rather than confirmatory.

\begin{table}[H]
\centering
\footnotesize
\begin{tabular}{lllccc}
\toprule
System & Experiment & Strategy pair & $p$ (final HV) & $p$ (AUC-HV) & $p$ (cum.\ MI) \\
\midrule
RHEA & Baseline & Adaptive-$\beta$ vs.\ Rule-Swap & $0.090$ & $0.006$\textsuperscript{**} & $<0.001$\textsuperscript{***} \\
 &  & Adaptive-$\beta$ vs.\ Fixed Mixed & $0.080$ & $<0.001$\textsuperscript{***} & $<0.001$\textsuperscript{***} \\
 &  & Adaptive-$\beta$ vs.\ qEHVI & $<0.001$\textsuperscript{***} & $<0.001$\textsuperscript{***} & $<0.001$\textsuperscript{***} \\
 &  & Rule-Swap vs.\ Fixed Mixed & $0.958$ & $<0.001$\textsuperscript{***} & $0.006$\textsuperscript{**} \\
 &  & Rule-Swap vs.\ qEHVI & $<0.001$\textsuperscript{***} & $<0.001$\textsuperscript{***} & $<0.001$\textsuperscript{***} \\
 &  & qEHVI vs.\ Fixed Mixed & $<0.001$\textsuperscript{***} & $<0.001$\textsuperscript{***} & $0.026$\textsuperscript{*} \\
\midrule
RHEA & Budget cut & Adaptive-$\beta$ vs.\ Rule-Swap & $0.059$ & $0.210$ & $0.011$\textsuperscript{*} \\
 &  & Adaptive-$\beta$ vs.\ Fixed Mixed & $0.075$ & $<0.001$\textsuperscript{***} & $0.182$ \\
 &  & Adaptive-$\beta$ vs.\ qEHVI & $0.007$\textsuperscript{**} & $0.017$\textsuperscript{*} & $<0.001$\textsuperscript{***} \\
 &  & Rule-Swap vs.\ Fixed Mixed & $0.275$ & $<0.001$\textsuperscript{***} & $0.210$ \\
 &  & Rule-Swap vs.\ qEHVI & $0.096$ & $0.001$\textsuperscript{**} & $<0.001$\textsuperscript{***} \\
 &  & qEHVI vs.\ Fixed Mixed & $0.012$\textsuperscript{*} & $<0.001$\textsuperscript{***} & $<0.001$\textsuperscript{***} \\
\midrule
RHEA & Time cut & Adaptive-$\beta$ vs.\ Rule-Swap & $0.173$ & $0.071$ & $0.017$\textsuperscript{*} \\
 &  & Adaptive-$\beta$ vs.\ Fixed Mixed & $0.937$ & $<0.001$\textsuperscript{***} & $0.007$\textsuperscript{**} \\
 &  & Adaptive-$\beta$ vs.\ qEHVI & $0.005$\textsuperscript{**} & $0.007$\textsuperscript{**} & $<0.001$\textsuperscript{***} \\
 &  & Rule-Swap vs.\ Fixed Mixed & $0.200$ & $<0.001$\textsuperscript{***} & $0.874$ \\
 &  & Rule-Swap vs.\ qEHVI & $0.059$ & $0.020$\textsuperscript{*} & $0.001$\textsuperscript{**} \\
 &  & qEHVI vs.\ Fixed Mixed & $0.001$\textsuperscript{**} & $<0.001$\textsuperscript{***} & $<0.001$\textsuperscript{***} \\
\midrule
RHEA & Dual & Adaptive-$\beta$ vs.\ Rule-Swap & $0.191$ & $0.692$ & $0.012$\textsuperscript{*} \\
 &  & Adaptive-$\beta$ vs.\ Fixed Mixed & $0.042$\textsuperscript{*} & $0.001$\textsuperscript{**} & $0.080$ \\
 &  & Adaptive-$\beta$ vs.\ qEHVI & $0.011$\textsuperscript{*} & $0.148$ & $<0.001$\textsuperscript{***} \\
 &  & Rule-Swap vs.\ Fixed Mixed & $0.230$ & $<0.001$\textsuperscript{***} & $0.596$ \\
 &  & Rule-Swap vs.\ qEHVI & $0.059$ & $0.002$\textsuperscript{**} & $0.032$\textsuperscript{*} \\
 &  & qEHVI vs.\ Fixed Mixed & $0.937$ & $<0.001$\textsuperscript{***} & $0.027$\textsuperscript{*} \\
\midrule
Fe-Co-Ni & Baseline & Adaptive-$\beta$ vs.\ Rule-Swap & $0.596$ & $0.833$ & $<0.001$\textsuperscript{***} \\
 &  & Adaptive-$\beta$ vs.\ Fixed Mixed & $0.634$ & $0.958$ & $<0.001$\textsuperscript{***} \\
 &  & Adaptive-$\beta$ vs.\ qEHVI & $0.672$ & $0.411$ & $<0.001$\textsuperscript{***} \\
 &  & Rule-Swap vs.\ Fixed Mixed & $0.711$ & $0.731$ & $0.010$\textsuperscript{**} \\
 &  & Rule-Swap vs.\ qEHVI & $0.653$ & $0.542$ & $<0.001$\textsuperscript{***} \\
 &  & qEHVI vs.\ Fixed Mixed & $0.542$ & $0.542$ & $<0.001$\textsuperscript{***} \\
\midrule
Fe-Co-Ni & Budget cut & Adaptive-$\beta$ vs.\ Rule-Swap & $0.312$ & $0.596$ & $<0.001$\textsuperscript{***} \\
 &  & Adaptive-$\beta$ vs.\ Fixed Mixed & $0.045$\textsuperscript{*} & $0.080$ & $<0.001$\textsuperscript{***} \\
 &  & Adaptive-$\beta$ vs.\ qEHVI & $0.711$ & $0.560$ & $<0.001$\textsuperscript{***} \\
 &  & Rule-Swap vs.\ Fixed Mixed & $0.791$ & $0.560$ & $0.071$ \\
 &  & Rule-Swap vs.\ qEHVI & $0.367$ & $0.191$ & $<0.001$\textsuperscript{***} \\
 &  & qEHVI vs.\ Fixed Mixed & $0.312$ & $0.026$\textsuperscript{*} & $0.037$\textsuperscript{*} \\
\midrule
Fe-Co-Ni & Time cut & Adaptive-$\beta$ vs.\ Rule-Swap & $0.411$ & $0.127$ & $0.011$\textsuperscript{*} \\
 &  & Adaptive-$\beta$ vs.\ Fixed Mixed & $0.474$ & $0.141$ & $<0.001$\textsuperscript{***} \\
 &  & Adaptive-$\beta$ vs.\ qEHVI & $0.711$ & $0.812$ & $<0.001$\textsuperscript{***} \\
 &  & Rule-Swap vs.\ Fixed Mixed & $0.080$ & $0.027$\textsuperscript{*} & $0.127$ \\
 &  & Rule-Swap vs.\ qEHVI & $0.164$ & $0.075$ & $0.001$\textsuperscript{**} \\
 &  & qEHVI vs.\ Fixed Mixed & $0.833$ & $0.141$ & $0.164$ \\
\midrule
Fe-Co-Ni & Dual & Adaptive-$\beta$ vs.\ Rule-Swap & $0.525$ & $0.164$ & $<0.001$\textsuperscript{***} \\
 &  & Adaptive-$\beta$ vs.\ Fixed Mixed & $0.615$ & $0.791$ & $<0.001$\textsuperscript{***} \\
 &  & Adaptive-$\beta$ vs.\ qEHVI & $0.474$ & $0.063$ & $<0.001$\textsuperscript{***} \\
 &  & Rule-Swap vs.\ Fixed Mixed & $0.367$ & $0.067$ & $0.011$\textsuperscript{*} \\
 &  & Rule-Swap vs.\ qEHVI & $0.711$ & $0.958$ & $<0.001$\textsuperscript{***} \\
 &  & qEHVI vs.\ Fixed Mixed & $0.833$ & $0.134$ & $0.113$ \\
\bottomrule
\end{tabular}

\caption{Pairwise strategy comparisons for final normalized hypervolume, AUC-HV,
         and cumulative mutual information ($n = 25$ seeds per strategy). Two-sided
         paired Wilcoxon signed-rank, uncorrected; markers denote
         \textsuperscript{*}$p < 0.05$, \textsuperscript{**}$p < 0.01$,
         \textsuperscript{***}$p < 0.001$.}
\label{tab:significance}
\end{table}
